\section{Initial machine design procedure and magnetic circuit sizing}

\subsection{Pole number ($2p$), synchronous speed ($n$), and stator slot count ($S_1, q_1$)}
\subsubsection{Pole pairs $p$ and synchronous speed}
Synchronous speed in r/min and r/s:
\[
n_s = \frac{60 f}{p}\quad[\mathrm{r/min}]
\quad\iff\quad
\boxed{p = \frac{60 f}{n_s}},
\qquad
\boxed{2p = \frac{120 f}{n_s}},
\qquad
\boxed{n = \frac{n_s}{60} = \frac{f}{p}\quad[\mathrm{r/s}]}.
\]

\subsubsection{Slots per pole per phase $q_1$ and total stator slots $S_1$}
Slots per pole per phase:
\[
q_1 = \frac{S_1}{2p\,m}\quad(m=3\text{ for three-phase}).
\]
Total number of stator slots:
\[
\boxed{S_1 = q_1\cdot(2p)\cdot m = 6\,p\,q_1}.
\]
\begin{itemize}
  \item Integer-slot winding: $q_1 \in \{1, 2, 3, 4, \dots\}$ (typically $q_1 = 2, 3, \text{ or } 4$).
  \item Fractional-slot winding: $q_1 = a + b/c$ (reduces tooth harmonics and cogging torque).
\end{itemize}
Stator slot pitch at bore circumference:
\[
\boxed{\tau_s = \frac{\pi D_g}{S_1}}\quad[\mathrm{m}].
\]

\subsection{Magnetic circuit dimensions and flux density limits ($B_t, B_y, b_t, b_y$)}
\subsubsection{Pole flux $\Phi_m$ and flux division in stator and rotor yokes}
Pole flux:
\[
\boxed{\Phi_m = \bar{B}\,\tau\,L_c = \bar{B}\left(\frac{\pi D_g}{2p}\right)L_c}\quad[\mathrm{Wb}].
\]
\textbf{Flux division rule in yokes (core backs):}
The magnetic flux entering the stator core from one pole divides equally into two halves ($\Phi_m/2$) travelling in opposite circumferential directions around the stator yoke.
The same half-division occurs in the rotor yoke.
\[
\boxed{\Phi_{\mathrm{yoke}} = \frac{\Phi_m}{2}}.
\]

\subsubsection{Stator tooth width $b_t$ and slot width $b_s$}
Number of teeth per pole is $\frac{S_1}{2p} = m q_1$.
Assuming all flux per pole passes through the stator teeth under that pole:
\[
\Phi_m = \left(\frac{S_1}{2p}\right)\cdot (b_t L_c)\cdot B_t
\quad\implies\quad
\boxed{b_t = \frac{\Phi_m}{(S_1/2p)\,B_t\,L_c} = \frac{\Phi_m}{m\,q_1\,B_t\,L_c}}\quad[\mathrm{m}],
\]
where $B_t$ is maximum allowable tooth flux density (typically $B_t \approx 1.5\text{--}1.8\,\mathrm{T}$ at the narrowest tooth section).
Slot width at bore:
\[
\boxed{b_s = \tau_s - b_t = \frac{\pi D_g}{S_1} - b_t}\quad[\mathrm{m}].
\]

\subsubsection{Stator yoke depth $b_y$ (core back)}
Since the stator yoke carries $\Phi_m/2$:
\[
\frac{\Phi_m}{2} = B_y\cdot(b_y L_c)
\quad\implies\quad
\boxed{b_y = \frac{\Phi_m/2}{B_y\,L_c} = \frac{\Phi_m}{2 B_y L_c}}\quad[\mathrm{m}],
\]
where $B_y$ is stator yoke flux density (typically $B_y \approx 1.1\text{--}1.5\,\mathrm{T}$).
Stator outer core diameter:
\[
\boxed{D_{\mathrm{out}} = D_g + 2 h_s + 2 b_y}\quad[\mathrm{m}].
\]

\subsubsection{Rotor pole body and pole shoe dimensions (salient-pole sync machine)}
Cross-sectional widths for rotor pole body and pole head:
\[
\boxed{b_{\mathrm{pole\_body}} = \frac{\Phi_m}{B_{\mathrm{pole\_body}}\,L_c}},
\qquad
\boxed{b_{\mathrm{pole\_head}} = \frac{\Phi_m}{B_{\mathrm{pole\_head}}\,L_c}},
\]
where typically $B_{\mathrm{pole\_body}} \approx 1.4\text{--}1.6\,\mathrm{T}$.

\subsection{Stator slot geometry, conductors per slot $z_s$, and slot depth $h_s$}
\subsubsection{Conductors per slot $z_s$ and conductor area $A_{\mathrm{cond}}$}
Turns per phase from phase EMF $E_{\mathrm{ph}} \approx V_{\mathrm{ph}}$:
\[
T_{\mathrm{ph}} = \frac{E_{\mathrm{ph}}}{4.44\,K_w\,\Phi_m\,f}.
\]
Conductors per slot:
\[
\boxed{z_s = \frac{2 T_{\mathrm{ph}}}{S_1/m} = \frac{2 m T_{\mathrm{ph}}}{S_1} = \frac{T_{\mathrm{ph}}}{p q_1}}.
\]
Conductor cross-sectional area with stator current density $J_s$ ($3\text{--}6\,\mathrm{A/mm^2}$):
\[
\boxed{A_{\mathrm{cond}} = \frac{I_c}{J_s} = \frac{I_{\mathrm{ph}}/a_p}{J_s}}\quad[\mathrm{mm^2}\text{ or }\mathrm{m^2}].
\]

\subsubsection{Slot gross area $A_{\mathrm{slot}}$ and slot depth $h_s$ with fill factor $k_{\mathrm{fill}}$}
Total conductor copper area per slot:
\[
A_{\mathrm{slot,cond}} = z_s A_{\mathrm{cond}}.
\]
Gross slot area with slot space factor $k_{\mathrm{fill}}$ ($\sim 0.35\text{--}0.45$ for round wire):
\[
\boxed{A_{\mathrm{slot}} = \frac{A_{\mathrm{slot,cond}}}{k_{\mathrm{fill}}} = \frac{z_s A_{\mathrm{cond}}}{k_{\mathrm{fill}}}}.
\]
Slot depth $h_s$ (parallel-sided slot):
\[
\boxed{h_s = \frac{A_{\mathrm{slot}}}{b_s} = \frac{z_s A_{\mathrm{cond}}}{k_{\mathrm{fill}} b_s}}\quad[\mathrm{m}].
\]
Slot aspect ratio $h_s / b_s$ is typically between 3 and 6 (excessively deep slots increase slot leakage inductance).

\subsection{Excitation MMF ($N_f I_f$) and rotor field winding sizing}
\subsubsection{Total magnetic circuit excitation MMF sum ($\sum H_k \ell_k$)}
Ampere's circuital law around one complete magnetic pole pair loop:
Total required MMF per pole:
\[
\boxed{(N_f I_f) = H_g g' + H_t h_t + H_{y,s} \ell_{y,s} + H_{\mathrm{pole\_head}} h_{\mathrm{pole\_head}} + H_{\mathrm{pole\_body}} h_{\mathrm{pole\_body}} + H_{y,r} \ell_{y,r}}\quad[\mathrm{A\text{-}turns}].
\]
\begin{itemize}
  \item Air gap represents the largest single MMF requirement: $F_g = H_g g' = \frac{B_g g'}{\mu_0} = \frac{B_g K_g g}{\mu_0}$.
  \item Iron path MMF drops: $F_k = H_k \ell_k$, where $H_k$ is found from the core steel $B$-$H$ curve at the operating flux density $B_k$.
\end{itemize}

\subsubsection{Rotor slot area $A_{\mathrm{slot,r}}$ from field current density $J_f$}
Field current density $J_f = \frac{I_f}{A_{\mathrm{cond,f}}}$ gives:
\[
A_{\mathrm{slot,r}} = N_f A_{\mathrm{cond,f}} = N_f \left(\frac{I_f}{J_f}\right) = \frac{N_f I_f}{J_f}
\quad\implies\quad
\boxed{A_{\mathrm{slot,r}} = \frac{(N_f I_f)}{J_f}}\quad[\mathrm{m^2}].
\]

\subsubsection{Field coil resistance $R_f$ and excitation voltage $U_f$}
Mean turn length $\ell_{\mathrm{mt}} \approx 2 L_c$. Using core length $L_c$ as axial length proxy:
\[
\boxed{R_f = \frac{N_f^2 L_c}{\sigma_{\mathrm{Cu}} A_{\mathrm{slot,r}}}}\quad[\Omega].
\]
Excitation terminal DC voltage $U_f$:
\[
\boxed{U_f = R_f I_f = \frac{(N_f I_f)\,N_f L_c}{\sigma_{\mathrm{Cu}} A_{\mathrm{slot,r}}}}\quad[\mathrm{V}].
\]
Solving for number of turns $N_f$ given excitation voltage $U_f$:
\[
\boxed{N_f = \frac{U_f\,\sigma_{\mathrm{Cu}} A_{\mathrm{slot,r}}}{(N_f I_f)\,L_c}},
\qquad
\boxed{I_f = \frac{(N_f I_f)}{N_f}}.
\]

\subsubsection{Field copper loss $P_{\mathrm{Cu,f}}$ and inverse slot area scaling ($P_{\mathrm{Cu,f}} \propto 1/A_{\mathrm{slot,r}}$)}
Field winding DC copper loss:
\[
\boxed{P_{\mathrm{Cu,f}} = R_f I_f^2 = \frac{N_f^2 L_c}{\sigma_{\mathrm{Cu}} A_{\mathrm{slot,r}}} I_f^2 = \frac{(N_f I_f)^2 L_c}{\sigma_{\mathrm{Cu}} A_{\mathrm{slot,r}}} \propto \frac{1}{A_{\mathrm{slot,r}}}}\quad[\mathrm{W}].
\]
\textbf{FUNDAMENTAL DESIGN THEOREM:}
For a given required excitation MMF $(N_f I_f)$, field copper loss depends ONLY on available rotor slot area $A_{\mathrm{slot,r}}$ and core length $L_c$. It is \textbf{independent} of the choice of turn number $N_f$ and current $I_f$!
To reduce field heating loss, one MUST increase rotor slot cross-sectional area $A_{\mathrm{slot,r}}$.

\subsection{Exam recipe: Synchronous machine initial design checklist and procedure}
\textbf{Use when:} given machine ratings ($S_{3\phi}, V_{\mathrm{ph}}, f, n, \bar{B}, ac$) and tasked with completing initial dimensions, slot counts, core sections, and rotor field winding.

\textbf{Steps:}
\begin{enumerate}
  \item \textbf{Speed, poles and stator slots:}
  Calculate $p = \frac{60 f}{n_s}$ [r/min] $\implies 2p = 2p$.
  Choose slots/pole/phase $q_1$ (e.g.\ $q_1 = 2$ or 3) $\implies$ total slots $S_1 = q_1 (2p) m$.
  Slot pitch $\tau_s = \frac{\pi D_g}{S_1}$.
  \item \textbf{Main dimensions and air gap:}
  Calculate output coefficient $G = 1.11 K_w \pi^2 \bar{B} ac$.
  Find volume proxy $D_g^2 L_c = \frac{S_{3\phi}}{G n}$ ($n = n_s/60$ in r/s).
  With square pole $\tau = L_c = \frac{\pi D_g}{2p}$, solve for $D_g$ and $L_c$. Choose air gap $g$.
  \item \textbf{Pole flux $\Phi_m$:}
  Compute $\Phi_m = \bar{B} \tau L_c = \bar{B}\left(\frac{\pi D_g}{2p}\right)L_c$.
  \item \textbf{Stator tooth width and slot width:}
  Choose tooth flux density $B_t \approx 1.5\text{--}1.8\,\mathrm{T}$.
  Calculate tooth width $b_t = \frac{\Phi_m}{(S_1/2p) B_t L_c}$ and slot width $b_s = \tau_s - b_t$.
  \item \textbf{Stator turns, conductors per slot, and slot depth:}
  Calculate turns $T_{\mathrm{ph}} = \frac{E_{\mathrm{ph}}}{4.44 K_w \Phi_m f}$.
  Conductors per slot $z_s = \frac{2 T_{\mathrm{ph}}}{S_1 / m}$ (round to even integer).
  Conductor area $A_{\mathrm{cond}} = \frac{I_{\mathrm{ph}}}{J_s}$.
  Slot depth $h_s = \frac{z_s A_{\mathrm{cond}}}{k_{\mathrm{fill}} b_s}$.
  \item \textbf{Stator yoke depth:}
  Choose yoke flux density $B_y \approx 1.2\text{--}1.5\,\mathrm{T}$.
  Calculate yoke depth $b_y = \frac{\Phi_m / 2}{B_y L_c}$ (\textbf{remember half flux!}).
  Outer diameter $D_{\mathrm{out}} = D_g + 2 h_s + 2 b_y$.
  \item \textbf{Rotor pole body and shoe:}
  Calculate widths $b_{\mathrm{pole\_body}} = \frac{\Phi_m}{B_{\mathrm{pole\_body}} L_c}$ and $b_{\mathrm{pole\_head}} = \frac{\Phi_m}{B_{\mathrm{pole\_head}} L_c}$.
  \item \textbf{Excitation MMF sum:}
  Sum MMFs: $(N_f I_f) = H_g g' + H_t h_t + H_{y,s} \ell_{y,s} + H_{\mathrm{pole}} h_{\mathrm{pole}} + H_{y,r} \ell_{y,r}$.
  \item \textbf{Rotor field slot and winding:}
  Given field current density $J_f$, find rotor slot area $A_{\mathrm{slot,r}} = \frac{N_f I_f}{J_f}$.
  Given DC excitation voltage $U_f$, solve $N_f = \frac{U_f \sigma_{\mathrm{Cu}} A_{\mathrm{slot,r}}}{(N_f I_f) L_c}$ and $I_f = \frac{N_f I_f}{N_f}$.
  Field loss $P_{\mathrm{Cu,f}} = \frac{(N_f I_f)^2 L_c}{\sigma_{\mathrm{Cu}} A_{\mathrm{slot,r}}}$.
\end{enumerate}

\textbf{Fast checks:}
\begin{itemize}
  \item $z_s$ must be an even integer for 2-layer stator windings.
  \item Stator yoke carries $\Phi_m/2$, NOT full $\Phi_m$.
  \item Rotor field copper loss is inversely proportional to rotor slot area: $P_{\mathrm{Cu,f}} \propto 1/A_{\mathrm{slot,r}}$.
\end{itemize}

\textbf{Common traps:}
\begin{itemize}
  \item Forgetting factor $1/2$ in yoke flux calculation ($b_y = \frac{\Phi_m/2}{B_y L_c}$).
  \item Mixing up $p$ (pole-pairs) and $2p$ (poles) when computing slot count $S_1 = q_1 (2p) m$.
  \item Using speed in r/min instead of r/s when calculating $G$ and volume proxy.
\end{itemize}
